\clearpage
\begin{explanationbox}
\textbf{\textcolor{CExplanation}{一、为什么会有这条结论？}}
\begin{center}\begin{tikzpicture}[font=\small,>=stealth,line join=round]
\draw[gray!40,rounded corners=3pt] (0,0) rectangle (14.2,8.3);\node at (7.1,7.700000000000001) {奇函数：把相加改成同一条曲线上的比较};
\begin{scope}[shift={(4.1,3.9499999999999997)},x=1.2cm,y=0.22999999999999998cm]\draw[->,gray] (-2.5,0)--(2.5,0) node[right]{$x$};\draw[->,gray] (0,-10)--(0,10) node[above]{$y$};\draw[blue,thick,solid,domain=-2.1:2.1,samples=180] plot(\x,{(\x)^3});\fill[orange!85!black] (1,1) circle(2.2pt);\fill[orange!85!black] (2,8) circle(2.2pt);\end{scope}\node[anchor=west,align=left,text width=5.0cm] at (8.4,4.0) {图取 $f(x)=x^3$\\[9pt]$a=2,\ b=-1$\\[9pt]$f(-b)=f(1)=1$\\[9pt]$f(a)+f(b)=8-1=7>0$\\[9pt]比较 $a=2>-b=1$};
\node[text width=12.6cm,align=center] at (7.1,.65) {图中标出的是比较用的 $a$ 和 $-b$；规律依赖奇性与严格单调性。};
\end{tikzpicture}\end{center}



想象老师只说：“这个函数是奇函数，在 $x\geqslant0$ 时越往右，图像越往上。”但老师不告诉你 $f(x)$ 到底等于什么。
如果题目让你解 $f(x-1)+f(2x)>0$，看起来两个函数值都算不出来。
\textbf{突破口是：我们不必算出它们，只需判断谁比谁大。}

先用一个初中生就能算的例子试试。取 $f(x)=x$，有
\[
 f(2)+f(-1)=2-1=1>0,\qquad 2+(-1)=1>0.
\]
再换一个不是直线的奇函数 $f(x)=x^3$，有
\[
 f(2)+f(-1)=8-1=7>0,\qquad 2+(-1)=1>0.
\]
\textbf{注意：}两次算出的函数值之和不是同一个数（$1$ 和 $7$），但是正负都由 $2+(-1)$ 决定。
下面说明：这不是碰巧，而是奇函数的对称性与严格递增共同保证的。

\medskip

\end{explanationbox}
\clearpage
\begin{explanationbox}
\textbf{\textcolor{CExplanation}{二、奇函数：为什么“加”能变成“比”？}}
\begin{center}\begin{tikzpicture}[font=\small,>=stealth,line join=round]
\draw[gray!40,rounded corners=3pt] (0,0) rectangle (14.2,8.3);\node at (7.1,7.700000000000001) {奇函数绕原点对称：左、中、右连成同向单调};
\begin{scope}[shift={(4.1,3.9499999999999997)},x=1.3636363636363635cm,y=0.32857142857142857cm]\draw[->,gray] (-2.2,0)--(2.2,0) node[right]{$x$};\draw[->,gray] (0,-7)--(0,7) node[above]{$y$};\draw[blue,thick,solid,domain=-1.8:1.8,samples=180] plot(\x,{(\x)^3});\fill[orange!85!black] (-1,-1) circle(2.2pt);\fill[orange!85!black] (0,0) circle(2.2pt);\fill[orange!85!black] (1,1) circle(2.2pt);\end{scope}\node[anchor=west,align=left,text width=5.0cm] at (8.4,4.0) {图取 $f(x)=x^3$\\[9pt]左侧 $f(u)<0$\\[9pt]原点 $f(0)=0$\\[9pt]右侧 $f(v)>0$\\[9pt]跨越 $0$ 仍有 $f(u)<f(v)$};
\node[text width=12.6cm,align=center] at (7.1,.65) {证明还要比较同在左侧的两点；用奇性将它们对应到右侧即可。};
\end{tikzpicture}\end{center}



奇函数最容易用上的等式是
\[
 f(-b)=-f(b).
\]
把要判断的式子一步步改写：
\[
\begin{aligned}
 f(a)+f(b)>0
 &\iff f(a)>-f(b)\\
 &\iff f(a)>f(-b).
\end{aligned}
\]
\textbf{原来我们根本不需要计算 $f(a)$、$f(b)$！}只需比较 $f(a)$ 与 $f(-b)$。

如果这条奇函数图像从左到右始终严格上升，那么：输入的数较大，输出的函数值一定较大。
因此
\[
 f(a)>f(-b)\iff a>-b\iff a+b>0.
\]

\textbf{关键疑问：题目只说右半边上升，为什么左半边也上升？}

奇函数满足 $f(0)=0$（因为 $f(0)=-f(0)$）。右半边严格上升，所以正数处函数值都大于 $0$；左半边是右半边的图像绕原点转半圈得到的，所以也从左往右上升，而且左侧函数值小于 $0$，右侧函数值大于 $0$。这就连成了整个区间的严格上升。

\textbf{再看一组正负数：}若 $a=-2,b=1$，那么 $a+b=-1<0$；取 $f(x)=x^3$，有 $f(-2)+f(1)=-8+1=-7<0$。规律与刚才一致。

\medskip

\end{explanationbox}
\clearpage
\begin{explanationbox}
\textbf{\textcolor{CExplanation}{三、偶函数：为什么要比绝对值？}}
\begin{center}\begin{tikzpicture}[font=\small,>=stealth,line join=round]
\draw[gray!40,rounded corners=3pt] (0,0) rectangle (14.2,8.3);\node at (7.1,7.700000000000001) {偶函数：先把负输入折到右边，再比距离};
\begin{scope}[shift={(4.1,2.003846153846154)},x=0.8108108108108107cm,y=0.3538461538461538cm]\draw[->,gray] (-3.7,0)--(3.7,0) node[right]{$x$};\draw[->,gray] (0,-1)--(0,12) node[above]{$y$};\draw[blue,thick,solid,domain=-3.2:3.2,samples=180] plot(\x,{(\x)^2});\fill[orange!85!black] (-3,9) circle(2.2pt);\fill[orange!85!black] (3,9) circle(2.2pt);\fill[orange!85!black] (2,4) circle(2.2pt);\draw[violet,dashed] (-3,9)--(3,9);\end{scope}\node[anchor=west,align=left,text width=5.0cm] at (8.4,4.0) {图取 $f(x)=x^2$\\[9pt]$a=-3,\ b=2$\\[9pt]$f(-3)=f(3)=9>4=f(2)$\\[9pt]比较 $|-3|=3>|2|=2$};
\node[text width=12.6cm,align=center] at (7.1,.65) {横坐标较小不等于函数值较小；真正比较的是到纵轴的水平距离。};
\end{tikzpicture}\end{center}



偶函数有 $f(-x)=f(x)$。例如 $f(x)=x^2$：
\[
 f(-3)=9,\quad f(3)=9,\quad f(1)=1.
\]
这里 $-3<1$，但 $f(-3)>f(1)$。\textbf{所以比 $-3$ 和 $1$ 的大小是错的。}

真正该比的是：两个点离 $y$ 轴（也就是 $x=0$）有多远。
$-3$ 离 $0$ 有 $3$ 格，$1$ 离 $0$ 有 $1$ 格；在右半边图像严格上升时，离得远的那个函数值更大。

每个数都可用绝对值表示“离 $0$ 的距离”，而偶函数左右一样高：
\[
 f(a)=f(|a|),\qquad f(b)=f(|b|).
\]
于是
\[
 f(a)>f(b)\iff f(|a|)>f(|b|)\iff |a|>|b|.
\]
\textbf{一句话：偶函数只认离原点的左右距离，不认左边的负号。}

\medskip

\end{explanationbox}
\clearpage
\begin{explanationbox}
\textbf{\textcolor{CExplanation}{四、两幅最简单的示意图}}
\begin{center}\begin{tikzpicture}[font=\small,>=stealth,line join=round]
\draw[gray!40,rounded corners=3pt] (0,0) rectangle (14.2,8.3);\node at (7.1,7.700000000000001) {偶函数：先把负输入折到右边，再比距离};
\begin{scope}[shift={(4.1,2.003846153846154)},x=0.8108108108108107cm,y=0.3538461538461538cm]\draw[->,gray] (-3.7,0)--(3.7,0) node[right]{$x$};\draw[->,gray] (0,-1)--(0,12) node[above]{$y$};\draw[blue,thick,solid,domain=-3.2:3.2,samples=180] plot(\x,{(\x)^2});\fill[orange!85!black] (-3,9) circle(2.2pt);\fill[orange!85!black] (3,9) circle(2.2pt);\fill[orange!85!black] (2,4) circle(2.2pt);\draw[violet,dashed] (-3,9)--(3,9);\end{scope}\node[anchor=west,align=left,text width=5.0cm] at (8.4,4.0) {图取 $f(x)=x^2$\\[9pt]$a=-3,\ b=2$\\[9pt]$f(-3)=f(3)=9>4=f(2)$\\[9pt]比较 $|-3|=3>|2|=2$};
\node[text width=12.6cm,align=center] at (7.1,.65) {横坐标较小不等于函数值较小；真正比较的是到纵轴的水平距离。};
\end{tikzpicture}\end{center}




\textit{图像只是帮助理解：奇函数可以不是直线，偶函数也可以不是抛物线。只要符合前面的条件，比较大小的规律就不会变。}

\medskip

\end{explanationbox}
\clearpage
\begin{explanationbox}
\textbf{\textcolor{CExplanation}{五、真正值得记住的三个发现}}
\begin{center}\begin{tikzpicture}[font=\small,>=stealth,line join=round]
\draw[gray!40,rounded corners=3pt] (0,0) rectangle (14.2,8.3);\node at (7.1,7.700000000000001) {偶函数的右半边递增：最小值在中间};
\begin{scope}[shift={(4.1,2.0951612903225807)},x=0.5263157894736842cm,y=0.14838709677419354cm]\draw[->,gray] (-5.7,0)--(5.7,0) node[right]{$x$};\draw[->,gray] (0,-3)--(0,28) node[above]{$y$};\draw[blue,thick,solid,domain=-5:5,samples=180] plot(\x,{(\x)^2});\fill[orange!85!black] (-4,16) circle(2.2pt);\fill[orange!85!black] (-2,4) circle(2.2pt);\fill[orange!85!black] (0,0) circle(2.2pt);\fill[orange!85!black] (2,4) circle(2.2pt);\end{scope}\node[anchor=west,align=left,text width=5.0cm] at (8.4,4.0) {用 $f(x)=x^2$ 示范\\[9pt]定义域 $[-5,5]$\\[9pt]$f(-4)=16,\ f(2)=4$\\[9pt]$f(0)=0$\\[9pt]$f(-2)=f(2)$\\[9pt]一般结论：$f(0)$ 最小};
\node[text width=12.6cm,align=center] at (7.1,.65) {最小值所在位置由对称和右半边递增决定；具体最小值不一定是0。};
\end{tikzpicture}\end{center}


\begin{enumerate}[leftmargin=*,itemsep=0.45em]
  \item \textbf{找“等于零”的位置往往比直接求值更快。}奇函数的 $f(a)+f(b)$ 在 $a=-b$ 时恰好等于 $0$。如果整个图像上升，$a$ 越过 $-b$，和就变正；如果下降，和就变负。
  \item \textbf{偶函数的最小值或最大值常在 $x=0$。}若右半边严格递增，则 $f(0)$ 是整个定义域的最小值；若右半边严格递减，则 $f(0)$ 是最大值。这是因为所有 $x$ 都满足 $|x|\geqslant0$。
  \item \textbf{一个式子里出现 $f$，不代表必须先求出 $f$。}有时题目给“奇偶性＋单调性”已经足够。高考中的抽象函数不等式，常常就是在考你能否绕开具体计算。
\end{enumerate}

\textbf{和已学知识联系：}这里用到的只有奇偶函数、数轴上的大小、绝对值的距离意义、单调性及一元不等式；其中奇函数 $f(0)=0$ 与 F016，奇偶性定义域要求与 F015，奇偶运算关系与 F017、F018 都能相互联系。
\end{explanationbox}
